\section{Conclusion}
\label{sec:conclusion}

We worked through the classical statistical-ML workflow on 14M randomized
advertising users and asked whether ranking by causal effect beats ranking
by response. For visits, yes at small budgets: the X-learner finds about
60\,\% more incremental visits at a 5\,\% budget, provided it has at least 3M
training rows. For conversions, no: no uplift learner beats ranking by
conversion probability at any budget, because the ad lifts conversion in
proportion to baseline propensity. Targeting still pays, since the top 10\,\%
by conversion probability captures about 89\,\% of the lift of treating everyone
at a tenth of the volume. Along the way, two corrections to naive analysis
proved larger than any modeling choice: adjusting for imperfect
randomization shrinks the visit effect by 28\,\%, and distinguishing assignment from
exposure changes the estimand by an order of magnitude.

\paragraph{Future work.} Several next steps are immediate.
\begin{enumerate}
  \item \textbf{Locate the imbalance mechanism.} Test whether the treatment
        imbalance comes from per-test pooling or label-dependent
        sub-sampling, for example by recovering test identity from the
        categorical features and re-running the estimators within tests.
  \item \textbf{Full-data causal forest.} Refit the causal forest and
        logistic regression on all 8.39M rows to check whether the
        conversion deficit of the forest is a sampling artifact.
  \item \textbf{Composite targets.} Uplift on a value-weighted outcome
        (visit-then-convert) might recover signal that neither outcome has
        alone.
  \item \textbf{Other splits and datasets.} Repeat the central comparison
        across resampled splits and on other public uplift benchmarks to test
        whether ``uplift wins on the dense outcome, not the sparse one'' holds
        beyond Criteo.
\end{enumerate}

We release the code, configuration, test suite, and all full-data metrics,
tables, and figures in the repository accompanying this paper.
